%%%%%%%%%%%%%%%%%%%%%%%%%%%%%%%%%%%%%%%%%%%%%%%%%%%%%%%%%%%%%%%%%%%%%%%%%%%%%
%%                                                                         %%
%%                   LaTeX Vorlage für die Studenten der                   %%
%%              Dualen Hochschule Baden-Württemberg Ravensburg             %%
%%                                                                         %%
%%  Die Vorlage orientiert sich an den Gestaltungsrichtlinien der DHBW RV  %%
%%                                                                         %%
%%                                                                         %%
%% Ersteller:        Markus Schutz (WI06)                                  %%
%% letzten Änderung: 26. November 2013                                     %%
%%                                                                         %%
%%                                                                         %%
%% Wichtiger Hinweis zur Verwendung dieser Vorlage:                        %%
%%                                                                         %%
%% Damit die Vorlage verwendet und die PDF richtig und vollständig erzeugt %%
%% werden kann bedarf des manuellen Aufrufs von makeindex. Diese können    %%
%% optional auch direkt im Editor eingerichtet werden.                     %%
%% TeXnicCenter (Ausgabe -> Ausgabeprofile definieren... (Alt + F7)        %%
%%               LaTeX => PDF -> Nachbearbeitung)                          %%
%%                                                                         %%
%% (1) Stichwortverzeichnis (falls verwendet)                              %%
%% makeindex -g -s styles\stichwortverzeichnis.ist vorlage                 %%
%% (2) Abkürzungsverzeichnis                                               %%
%% makeindex vorlage.nlo -s cl.ist -o vorlage.nls                     %%
%% (3) Glossar (falls verwendet)                                           %%
%% makeindex -s vorlage.ist -t vorlage.glg -o vorlage.gls vorlage.glo      %%
%%                                                                         %%
%% Bei der Erzeugung der PDF Datei (Anwendung des LaTeX => PDF Ausgabe-    %%
%% profiles) werden das Abkürzungsverzeichnis, das Stichwortverzeichnis    %%
%% und das Glossar jetzt richtig erzeugt. In Verbindung mit TeXnicCenter   %%
%% kann es beim automatisierten Aufruf von makeindex zu Probleme kommen.   %%
%% Ein manueller Aufruf funktioniert dagegen immer.                        %%
%%                                                                         %%
%% Wichtiger Hinweis:                                                      %%
%%                                                                         %%
%% Keine Änderungen an den Dateien im Verzeichnis "pages" vornehmen. Für   %%
%% die Arbeit beziehen sich alle Änderungen auf diese Datei und die        %%
%% Dateien im Verzeichnis "chapter".                                       %%
%% Für die Erstellung des Literaturverzeichnises empfiehlt sich die Ver-   %%
%% wendung von JabRef (http://jabref.sourceforge.net). Die Datei ist unter %%
%% dem Namen literatur.bib im Verzeichnis "literatur" zu speichern.        %%
%%                                                                         %%
%% Zur sinnvollen Nutzung dieser Vorlage empfiehlt es sich, die Dokus zu   %%
%% den eingebundenen Paketen durchzulesen. Sie sind im doc-Verzeichnis der %%
%% MiKTeX-Installation zu finden.                                          %%
%%                                                                         %%
%% Enthaltene Titelblätter:                                                %%
%%   - Seminararbeit                                                       %%
%%   - Projektarbeit                                                       %%
%%   - Bachelorarbeit                                                      %%
%%                                                                         %%
%%%%%%%%%%%%%%%%%%%%%%%%%%%%%%%%%%%%%%%%%%%%%%%%%%%%%%%%%%%%%%%%%%%%%%%%%%%%%

\documentclass[a4paper,12pt]{article}                                         % Schriftgröße, Layout, Papierformat, Art des Dokumentes
\usepackage[left=3cm,right=2cm,top=2cm,bottom=2cm,includehead]{geometry}      % Einstellungen der Seitenränder
\usepackage{ngerman}                                                          % neue Rechtschreibung
\usepackage[main=ngerman,english]{babel} 
\usepackage{lmodern}                                                  % deutsche Silbentrennung
\usepackage[utf8]{inputenc} 
%\usepackage[ansinew]{inputenc}                                                  % Umlaute
\usepackage[T1]{fontenc}
\usepackage{textcomp}
\usepackage[hyperfootnotes=false]{hyperref}                                   % pfd-Output [Fußnoten nicht verlinken]
\usepackage[nottoc]{tocbibind}                                                % Inhaltsverzeichniserweiterung (Inhaltsverzeichnis selbst ausblenden)
\usepackage{makeidx}                                                          % Index
\usepackage[intoc]{nomencl}                                                   % Abkürzungsverzeichnis
\usepackage{fancyhdr}                                                         % Fancy Header
\usepackage[round]{natbib}                                                    % Zitate (Erweiterung für Literaturverzeichnis)
\usepackage{amsmath}                                                          % Zurücksetzen der Tabellen- und Abbildungsnummerierung je Sektion
\usepackage[labelfont=bf,aboveskip=1mm]{caption}                              % Bild- und Tabellenunterschrift (fett)
\usepackage{setspace}                                                         % Zeilenabstand (vor footmisc laden!)
\usepackage[bottom,multiple,hang,marginal]{footmisc}                          % Fußnoten [Ausrichtung unten, Trennung durch Seperator bei mehreren Fußnoten]
\usepackage{graphicx}                                                         % Grafiken
\usepackage{tabularx}                                                         % erweiterte Tabellen
\usepackage{longtable}                                                        % mehrseitige Tabellen
\usepackage{color}                                                            % Farben
\usepackage{enumitem}                                                         % Befehl setlist (Zeilenabstand für itemize Umgebung auf 1 setzen)
\usepackage[formats]{listings}                                                         % Quelltexte
\usepackage{zref}                                                             % Verweise (Anhangsverweise)
\usepackage[toc,style=altlist,translate=false]{glossaries}                    % Glossar (nach hyperref, 
%\usepackage{glossaries-babel}                                                 % Glossar: Übersetzung im TOC
\usepackage{url}
\usepackage[autostyle]{csquotes}
\usepackage{framed}
\usepackage{lipsum}
\usepackage[final]{pdfpages}
\usepackage{float}															% Platzierung der Bilder
\definecolor{shadecolor}{gray}{0.9}
\definecolor{gray}{rgb}{0.4,0.4,0.4}
\definecolor{darkblue}{rgb}{0.0,0.0,0.6}
\definecolor{cyan}{rgb}{0.0,0.6,0.6}




%%%%%%%%%%%%%%%%%%%%%%%%%%%%%%%%%%%%%%%%%%%%%%%%%%%%%%%%%%%%%%%%%%%%%%%%%%%%%
%%                                                                         %%
%% \/   \/      Bitte hier die Änderungen zur Arbeit vornehmen     \/   \/ %%
%%                                                                         %%
%%%%%%%%%%%%%%%%%%%%%%%%%%%%%%%%%%%%%%%%%%%%%%%%%%%%%%%%%%%%%%%%%%%%%%%%%%%%%

%%%%%%%%%%%%%%%%%%%%%%% Definitionen bzgl. der Arbeit %%%%%%%%%%%%%%%%%%%%%%%
\def\myType{3}           % [0=Seminararbeit|1=Projektarbeit|2=Bachelorarbeit|3=Studienarbeit]

\def\myTopic{Prototypische Implementierung eines Dozentenfinders mittels Bluetooth-Technologie}
\def\mySubTopic{ zur Machbarkeits-Evaluierung des “DeSearch”-Projekts }
\def\myAutor{Philipp Stehle und Lina Hirschoff}
\def\myMatNr{1234567 und 4125931}
\def\myCompany{SAP SE}
\def\myCompanyAddressStreet{Dornierstraße 3}
\def\myCompanyAddressCity{88677 Markdorf}
\def\myProf{Andreas Judt}
\def\myEndDate{15.07.2016}
\def\myEditDate{KW 40/2015 - KW 28/2016}

%%%%%%%%%%%%%%%%%%%% Folgende Angaben für: Seminararbeit %%%%%%%%%%%%%%%%%%%%
\def\myVorlesung{Name der Vorlesung}

%%%%%%%%%%%%%%%%%%%% Folgende Angaben für: Projektarbeit %%%%%%%%%%%%%%%%%%%%
\def\myProjNumber{4}         % [1|2]
\def\myPraxPhase{3}                       % [1|2|3]

%%%%%%%%%%%%%%%%%%%%%%%%%%%%%%%%%%%%%%%%%%%%%%%%%%%%%%%%%%%%%%%%%%%%%%%%%%%%%
%%                                                                         %%
%% /\   /\         Ab hier keine Änderungen mehr vornehmen         /\   /\ %%
%%                                                                         %%
%%%%%%%%%%%%%%%%%%%%%%%%%%%%%%%%%%%%%%%%%%%%%%%%%%%%%%%%%%%%%%%%%%%%%%%%%%%%%

%%%%%%%%%%%%%%%%%%%%%%%% Eigene Farbwerte definieren %%%%%%%%%%%%%%%%%%%%%%%%
\definecolor{boxgray}{gray}{0.9}         % Hintergrundfarbe für Zitatboxen
\definecolor{commentgray}{gray}{0.5}     % Grau für Kommentare in Quelltexten
\definecolor{darkgreen}{rgb}{0,.5,0}     % Grün für Strings in Quelltexten
\definecolor{purple}{rgb}{0.44, 0.16, 0.39} %Lila für JS-Keywords in Quelltexten


%%%%%%%%%%%%%%%%%%%%%%%% Eigene Kommandos definieren %%%%%%%%%%%%%%%%%%%%%%%%
% Definition von \gqq{#1: text}: Text in Anführungszeichen
\newcommand{\gqq}[1]{\glqq #1\grqq}

% Definition von \footref{#1: label}
% Verweis auf bereits existierende Fußnoten mittels
\providecommand*{\footref}[1]{
	\begingroup
		\unrestored@protected@xdef\@thefnmark{\ref{#1}}
	\endgroup
\@footnotemark}

% Definition von \mypageref{#1: label}
% Kombination aus \ref{#1: label} und \pageref{#1: label}
\newcommand{\mypageref}[1]{\ref{#1} auf Seite \pageref{#1}}

% Definition von \myboxquote{#1: text}
% grau hinterlegte Quotation-Umgebung (für Zitate)
\newcommand{\myboxquote}[1]{
	\begin{quotation}
		\colorbox{boxgray}{\parbox{0.78\textwidth}{#1}}
	\end{quotation}
	\vspace*{1mm}
}

\makeatletter
\zref@newprop*{appsec}{}
\zref@addprop{main}{appsec}

% Definition von \applabel{#1: label}{#2: text}
% von \appsec{#1: text}{#2: label} zur Erzeugung des Labels verwendet)
\def\applabel#1#2{%
	\zref@setcurrent{appsec}{#2}%   
	\zref@wrapper@immediate{\zref@label{#1}}%
}

% Definition von \appref{#1: label}
% anstelle \ref{#1: label} zum referenzieren von Anhängen verwenden)
\def\appref#1{%
	\hyperref[#1]{\zref@extract{#1}{appsec}}%
}
\makeatother

% Definition von \appsection{#1: text}{#2: label}
% Ersetzt \section{#1: text} und \label{#2: label} für Anhänge)
\newcommand{\theappsection}[1]{Anhang \Alph{section}:~\protect #1}
\newcommand{\appsection}[2]{
	\addtocounter{section}{1}
	\phantomsection
	\addcontentsline{toc}{section}{\theappsection{#1}}
	\markboth{\theappsection{#1}}{}

	\section*{\theappsection{#1}}
	\applabel{#2}{Anhang \Alph{section}}
	\label{#2}
}

%%%%%%%%%%%%% Index, Abkürzungsverzeichnis und Glossar erstellen %%%%%%%%%%%%
\makeindex
\makenomenclature
\makeglossaries

% Festlegung der Art der Zitierung (Havardmethode: Abkuerzung Autor + Jahr) %
\bibliographystyle{dinat}

%%%%%%%%%%%%%%%%%%%%%%%%%%%%%%% PDF-Optionen %%%%%%%%%%%%%%%%%%%%%%%%%%%%%%%%
\hypersetup{
	bookmarksopen=false,
	bookmarksnumbered=true,
	bookmarksopenlevel=0,
	pdftitle=\myTopic,
	pdfsubject=\myTopic,
	pdfauthor=\myAutor,
	pdfborder=0 0 0,
	pdfstartview=Fit,
	pdfpagelayout=SinglePage
}

%%%%%%%%%%%%%%%%%%%%%%%%%%%% Kopf- und Fußzeile %%%%%%%%%%%%%%%%%%%%%%%%%%%%%
\pagestyle{fancy}
\fancyhf{}
\fancyhead[R]{\thepage}                         % Kopfzeile rechts bzw. außen
\renewcommand{\headrulewidth}{0.5pt}            % Kopfzeile rechts bzw. außen

%%%%%%%%%%%%%%%%%%%%%%%%% Layout und Beschriftungen %%%%%%%%%%%%%%%%%%%%%%%%%
\onehalfspacing                % Zeilenabstand: 1.5 (Standard: 1.2)
\setlist{noitemsep}            % Zeilenabstand für items auf 1 setzen

\addto\captionsngerman{        % Tabllen- und Abbildungsunterschriften ändern
  \renewcommand{\figurename}{Abb.}
  \renewcommand{\tablename}{Tab.}
}
\numberwithin{table}{section}                               % Tabellennummerierung je Sektion zurücksetzen
\numberwithin{figure}{section}                              % Abbildungsnummerierung je Sektion zurücksetzen
\renewcommand{\thetable}{\arabic{section}.\arabic{table}}   % Tabellennummerierung mit Section
\renewcommand{\thefigure}{\arabic{section}.\arabic{figure}} % Abbildungsnummerierung mit Section
\renewcommand{\thefootnote}{\arabic{footnote}}              % Sektionsbezeichnung von Fußnoten entfernen

\renewcommand{\multfootsep}{, }                             % Mehrere Fußnoten durch ", " trennen

%%%%%%%%%%%%%%%%%%%%%%%%%%%%%%% Listingstyle %%%%%%%%%%%%%%%%%%%%%%%%%%%%%%%%
\lstset{
%	basicstyle=\ttfamily\small
%	basicstyle=\scriptsize\ttfamily\bfseries,
%	basicstyle=\fontfamily{pcr}\selectfont\scriptsize,
	%basicstyle=\itshape\small,
 basicstyle=\footnotesize\ttfamily,
	commentstyle=\color{commentgray}\textit,
	showstringspaces=false,
	stringstyle=\color{darkgreen},
	keywordstyle=\color{blue},
%	numbers=left,
%	numberstyle=\tiny,
%	stepnumber=1,
%	numbersep=10pt,
	tabsize=2,
	showspaces=false,
	keepspaces=true
	showtabs=false,
	fontadjust=true,
	frame=single,
%	backgroundcolor=\color{boxgray},
	captionpos=b,
	linewidth=1.0\linewidth,
	xleftmargin=0\linewidth,
	breaklines=true,
	aboveskip=16pt
}
\lstdefinelanguage{XML}
{
  morestring=[b]",
  morestring=[s]{>}{<},
  morecomment=[s]{<?}{?>},
  stringstyle=\color{black},
  identifierstyle=\color{darkblue},
  keywordstyle=\color{cyan},
  morekeywords={xmlns,version,type}% list your attributes here
}
\lstdefinelanguage{JavaScript}{
  keywords={break, case, catch, continue, debugger, default, delete, do, else, false, finally, for, function, if, in, instanceof, new, null, return, switch, this, throw, true, try, typeof, var, void, while, with},
  morecomment=[l]{//},
  morecomment=[s]{/*}{*/},
  morestring=[b]',
  morestring=[b]",
  ndkeywords={class, export, boolean, throw, implements, import, this},
  keywordstyle=\color{blue}\bfseries,
  ndkeywordstyle=\color{darkgray}\bfseries,
  identifierstyle=\color{black},
  commentstyle=\color{purple}\ttfamily,
  stringstyle=\color{darkgreen}\ttfamily,
  sensitive=true
}

        
%%%%%%%%%%%%%%%%%%%%%%%%%%%%%%%%%%%%%%%%%%%%%%%%%%%%%%%%%%%%%%%%%%%%%%%%%%%%%
%%                                                                         %%
%% \/   \/      Bitte hier die Änderungen zur Arbeit vornehmen     \/   \/ %%
%%                                                                         %%
%%%%%%%%%%%%%%%%%%%%%%%%%%%%%%%%%%%%%%%%%%%%%%%%%%%%%%%%%%%%%%%%%%%%%%%%%%%%%

%Seiten und Kapitel einbinden
\begin{document}
	\pagenumbering{Roman}
	% Das Titelblatt wird automatisch ausgewählt. Keine Änderung hier
	\ifcase\myType
		\begin{titlepage}
	\begin{center}
		\vspace*{2cm}
		\LARGE\bf\myTopic\\
		\Large\rm\mySubTopic\\
		\vspace*{3cm}
		\bf Seminararbeit zur Vorlesung\\
		\myVorlesung\\
		\normalsize\rm
		\vspace*{1cm}
		für die\\
		Prüfung zum Bachelor of Engineering\\
		\vspace*{1cm}
		an der Fakultät für Technik\\
		im Studiengang Informationstechnik\\
		\vspace*{1cm}
		an der\\
		DHBW Ravensburg\\
		Campus Friedrichshafen
		\vfill
	\end{center}
	\begin{tabular}{ll}
		Verfasser:&\myAutor\\
		Ausbildungsbetrieb:&\myCompany\\
		Anschrift:&\myCompanyAddressStreet\\
		&\myCompanyAddressCity\\
	%	Wiss. Betreuer:&\myProf\\
		Abgabedatum:&\myEndDate\\
	\end{tabular}
\end{titlepage}
\newpage
\setcounter{page}{2}

	\or
		\begin{titlepage}

\begin{table}
\begin{tabular}{p{0.4\textwidth}p{0.4\textwidth}}
\includegraphics[width=4cm]{images/sap_logo.png} &
\hfill \includegraphics[width=4cm]{images/dhbw_logo.jpg} \\
\end{tabular}
\end{table}

	\begin{center}
		\vspace*{1cm}
		\LARGE\bf\myTopic\\
		%\Large\rm\mySubTopic\\
		\vspace*{1cm}
		\bf Projektarbeit\\
		
		\normalsize\rm
		 im \myPraxPhase. Studienjahr \\
		\vspace*{1cm}
		
		Duale Hochschule Baden-Württemberg\\
		\vspace*{0.5cm}
		TIT13 Informationstechnik\\
		\vspace*{1cm}
		von\\
		 \myAutor\\
		\vspace*{1cm}
		
		\vfill
	\end{center}
	\begin{tabular}{ll}
		Abgabe:&\myEndDate\\
		Bearbeitungszeitraum:&\myEditDate\\
		Matrikelnummber:&\myMatNr\\
		Partnerunternehmen:&\myCompany\\
		Addresse:&\myCompanyAddressStreet\\
					&\myCompanyAddressCity\\
		\\
		Betreuer:&\myProf\\
		
	\end{tabular}
	\newline
	\vspace*{1cm}
%	\newline
%	\begin{tabularx}{\textwidth}{l@{\extracolsep\fill}r}
%	  Unterschrift des verantwortlichen Ausbilders&\\
%	  (oder des Personalverantwortlichen)&\rule{6cm}{0.3mm}\\
%	\end{tabularx}
\end{titlepage}
\newpage
\setcounter{page}{2}

	\or
		\begin{titlepage}
	\begin{center}
		\vspace*{2cm}
		\LARGE\bf\myTopic\\
		\Large\rm\mySubTopic\\
		\vspace*{3cm}
		\bf Bachelorarbeit\\
		\normalsize\rm
		\vspace*{1cm}
		für die\\
		Prüfung zum Bachelor of Science\\
		\vspace*{1cm}
		an der Fakultät für Wirtschaft\\
		im Studiengang Wirtschaftsinformatik\\
		\vspace*{1cm}
		an der\\
		DHBW Ravensburg
		\vfill
	\end{center}
	\begin{tabular}{ll}
		Verfasser:&\myAutor\\
		Ausbildungsbetrieb:&\myCompany\\
		Anschrift:&\myCompanyAddressStreet\\
		&\myCompanyAddressCity\\
		Wiss. Betreuer:&\myProf\\
		Abgabedatum:&\myEndDate\\
	\end{tabular}
\end{titlepage}
\newpage
\setcounter{page}{2}

	\or
		\begin{titlepage}

\begin{table}
\begin{tabular}{p{0.4\textwidth}p{0.4\textwidth}}
%\includegraphics[width=4cm]{images/sap_logo.png} &
\hfill \includegraphics[width=6cm]{images/dhbw_logo.jpg} \\
\end{tabular}
\end{table}

	\begin{center}
		\vspace*{1cm}
		\LARGE\bf\myTopic\\
		\Large\rm\mySubTopic\\
		\vspace*{1cm}
		\bf Studienarbeit\\
		
		\normalsize\rm
		 im 3. Studienjahr \\
		\vspace*{1cm}
		
		Duale Hochschule Baden-Württemberg\\
		\vspace*{0.5cm}
		TIT13 Informationstechnik\\
		\vspace*{1cm}
		von\\
		 \myAutor\\
		\vspace*{1cm}
		
		\vfill
	\end{center}
	\begin{tabular}{ll}
		Abgabe:&\myEndDate\\
		Bearbeitungszeitraum:&\myEditDate\\
		Matrikelnummer:&\myMatNr\\
		
		Betreuer:&\myProf\\
		
	\end{tabular}
	\newline
	\vspace*{1cm}
%	\newline
%	\begin{tabularx}{\textwidth}{l@{\extracolsep\fill}r}
%	  Unterschrift des verantwortlichen Ausbilders&\\
%	  (oder des Personalverantwortlichen)&\rule{6cm}{0.3mm}\\
%	\end{tabularx}
\end{titlepage}
\newpage
\setcounter{page}{2}

	\else
	\fi
	
	\pagestyle{fancy}
	\thispagestyle{empty}
\addcontentsline{toc}{section}{Selbständigkeitserklärung}
\begin{center}
	\vspace*{2cm}
	\Huge\bf Selbständigkeitserklärung\\
	\vspace*{2cm}
	\normalsize\rm
	gemäß Ziffer 1.2.3 der Anlage 1 zu §§ 3, 4, und 5  der „Studien- und Prüfungsordnung DHBW Technik“ vom 29. September 2015.\\
	\vspace*{1cm}
	\normalsize\rm
	Ich versichere hiermit, dass ich meine \ifcase\myType Seminararbeit \or Projektarbeit \or Bachelorarbeit \or Studienarbeit \else\fi mit dem Thema: \\
	\vspace*{1cm}
	\Large\bf\myTopic\\
	\Large\rm\mySubTopic\\
	\vspace*{1cm}
	\normalsize\rm
	selbständig verfasst und keine anderen als die angegebenen\\Quellen und Hilfsmittel benutzt habe. Ich versichere zudem, dass die eingereichte elektronische Fassung mit der gedruckten Fassung übereinstimmt.\\
	\vspace*{4cm}
	\begin{tabularx}{\textwidth}{l@{\extracolsep\fill}r}
  	\rule{7cm}{0.3mm}\\
	\end{tabularx}
	\begin{tabularx}{\textwidth}{*{2}{>{\arraybackslash}X}}
	  Ort, Datum\\
	\end{tabularx}
	
	\vfill
	\begin{tabularx}{\textwidth}{l@{\extracolsep\fill}r}
		\rule{7cm}{0.3mm}&\rule{7.55cm}{0.3mm}\\
	\end{tabularx}
	\begin{tabularx}{\textwidth}{*{2}{>{\arraybackslash}X}}
		Unterschrift&Unterschrift\\
	\end{tabularx}
\end{center}

%	\begin{titlepage}
	\begin{center}
		\vspace*{1cm}
		\Huge\bf Sperrvermerk\\
		\vspace*{2cm}
		\normalsize\rm
		\begin{quotation}
			\parbox{0.8\textwidth}{Der Inhalt dieser \ifcase\myType Seminararbeit \or Projektarbeit \or Bachelorarbeit\else\fi darf weder als Ganzes noch in Auszügen Personen außerhalb des Prüfungsprozesses und des Evaluationsverfahrens zugänglich gemacht werden, sofern keine anders lautende Genehmigung der Ausbildungsstätte vorliegt. Ausnahmen bedürfen der schriftlichen Genehmigung der Firma \myCompany, \myCompanyAddressStreet, \myCompanyAddressCity.}
		\end{quotation}
		\vspace*{1cm}
%		\begin{quotation}
%		  \parbox{0.8\textwidth}{
%		  \begin{tabularx}{0.78\textwidth}{l@{\extracolsep\fill}l}
%				\rule{4cm}{0.3mm}&\rule{4cm}{0.3mm}\\
%	    	Ort, Datum&Unterschrift
%			\end{tabularx}}
%		\end{quotation}
	\end{center}
\end{titlepage}
\newpage
\setcounter{page}{3}

%	\include{00_signaturen}
	\section*{Abstract}
Deutsch
\begin{center}
	\rule{0.3\textwidth}{0.4mm}
\end{center}
Englisch
	\tableofcontents
\newpage

	\renewcommand{\nomname}{Abkürzungsverzeichnis}
\setlength{\nomlabelwidth}{.25\hsize}
\renewcommand{\nomlabel}[1]{#1 \dotfill}
\setlength{\nomitemsep}{-\parsep}
\printnomenclature
%\newpage


	\renewcommand{\glossaryname}{Glossar}
\printglossaries
	\listoffigures
%\newpage

	%\listoftables
\newpage

	%Listingnummering je Sektion zurücksetzen
\numberwithin{lstlisting}{section}
%Listingnummerierung mit Section
\renewcommand{\thelstlisting}{\arabic{section}.\arabic{lstlisting}}
%Listingsverzeichnis in das Inhaltsverzeichnis aufnehmen.
\renewcommand{\lstlistingname}{Listingsverzeichnis}
\renewcommand{\lstlistoflistings}{\begingroup
\tocchapter
\tocfile{\lstlistingname}{lol}
\endgroup}

\lstlistoflistings
\newpage

\renewcommand{\lstlistingname}{Listing}


	% Kapitel
	\pagestyle{fancy}
	\fancyhead[L]{\nouppercase{\leftmark}}                               % Kopfzeile links bzw. innen
	\pagenumbering{arabic}
	

	
	
	% Abschluss
	\bibliography{literatur/Studienarbeit_literatur}
\newpage

	%\renewcommand{\indexname}{Stichwortverzeichnis}
\printindex
\newpage

		
	% Anhang
	\renewcommand{\thetable}{\Alph{section}.\arabic{table}}              % Tabellennummerierung mit Section
	\renewcommand{\thefigure}{\Alph{section}.\arabic{figure}}            % Abbildungsnummerierung mit Section
	\renewcommand{\thelstlisting}{\Alph{section}.\arabic{lstlisting}}    % Listingsnummerierung mit Section
	
	\begin{appendix}
%	\include{chapter/30_appendix}
	\end{appendix}
	
	
\end{document}

%%%%%%%%%%%%%%%%%%%%%%%%%%%%%%%%%%%%%%%%%%%%%%%%%%%%%%%%%%%%%%%%%%%%%%%%%%%%%
%%                                                                         %%
%% /\   /\         Ab hier keine Änderungen mehr vornehmen         /\   /\ %%
%%                                                                         %%
%%%%%%%%%%%%%%%%%%%%%%%%%%%%%%%%%%%%%%%%%%%%%%%%%%%%%%%%%%%%%%%%%%%%%%%%%%%%%